\section{Lattice calculation of PDF}\label{sec:numerical}

\subsection{Lattice Matrix Elements}

In this subsection,  we give  the results of a Lattice-QCD calculation using clover valence fermions on  {the CLS $32^3\times96$ 2+1 flavor clover fermion ensemble H102 with Lattice spacing $a=0.086$~fm, pion mass $M_\pi=$ 356~MeV and box size $L\approx 2.7$~fm ($M_\pi L\approx 4.9$)~\cite{Bruno:2014jqa}. We use  $\kappa_l=0.136865$ and $C_{SW}=1.98625$ for the valence clover fermion. We apply APE smearing~\cite{Hasenfratz:2001hp} with size=2.5$a$ twice in the source/sink smearing and also in the quasi-PDF operator $O_\Gamma$, but not in the fermion propagators.}

\begin{table}[htbp]
\begin{center}
\caption{\label{table:p_modes}
Momentum modes   used in the NPR analysis. The four digits (in units of $2\pi/L$) in   brackets correspond to   three spatial momentum components and the energy  component. The $z$ direction can be selected by  setting   the Wilson link along any of the three spatial directions.  Thus these momentum modes approximately cover  three choices  of $\mu_R=\sqrt{-(p^R)^2}$ and several sets of $p^R_z= 2\pi i/L(i=0,1,2,...)$.  }
\begin{tabular}{c|ccc}
\hline 
$a^2\mu_R^2$ & Momentum modes $p^R$ \\
$ [1.109,1.118]$ &(5,2,0,0)   (4,3,1,5)   \\
$ [1.957,1.966]$ &(5,5,0,3)  (6,3,2,4)   (5,4,1,9) \\
$ [2.814,2.832]$ &(6,5,1,10)   (7,4,2,6)   (6,3,0,16)\\
\hline
\end{tabular}
\end{center}
\end{table}


First of all, we  will explore   the nonperturbative renormalization   in  RI/MOM scheme.  Following  Ref.~\cite{Chen:2017mzz} , we use Landau gauge fixed wall sources (while limiting the source in the  time slice range $t\in[32, 64]$ to avoid the boundary effect from   open boundary condition at $t$=0), and generate the propagators with the momentum modes listed in Table.~\ref{table:p_modes}.  The four digits in   brackets correspond to   three spatial component and the energy  component  in units of $2\pi/L$. The $z$ direction can be selected by  setting   the Wilson link along any of the three spatial directions.   Thus these results approximately  cover three values of  {$\mu_R=\sqrt{-(p^R)^2}$ (2.4, 3.2 and 3.9~GeV, corresponding to $a^2\mu_R^2$=1.1, 2.0 and 2.8), and     $p^R_z=\{0,1,2,...,\}*2\pi /L$} up to the upper limit $p^R_z<\mu_R$. Note that in deriving these momentum modes we have required  the spatial components of a given momentum mode different with  each other, and adjusted  $p_t^R$ to ensure  $\mu_R$ invariant (within 2\%). These  choices allow us to explore  the dependence on each component of $p$, but one should be cautious that  the results may suffer from sizable  discretization errors since the normal constraint $\frac{\sum_{\mu}a^4p_{\mu}^4}{(\sum_{\mu}a^2p_{\mu}^2)^2}<0.3$ is not respected. These  discretization errors will be investigated in the future.


\begin{figure}[htbp]
\includegraphics[width=.48\textwidth]{images/compare_z6.pdf}
\includegraphics[width=.48\textwidth]{images/compare_delta_z6.pdf}
\caption{The NPR  $Z=Z_{mp}$ (top) and $\delta Z=Z_{\slashed{p}}-Z_{mp}$ (bottom) at $z=6a$ ($\approx 0.5$~fm) as a function of $p_z^R/\mu_R=1/\sqrt{r}$, with various $\mu_R$.   At $p_z^R=0$, $Z$ is real and $\delta Z/Z$ is less than 5\%. } \label{fig:pz_dependence}
\end{figure}

As shown in Eq.~(\ref{eq:matching_coeff}), the one-loop matching formula   primarily depends on the combination $r=\mu^2_R/(p_z^R)^2$ but  is independent of $p^R_t$. In Fig.~\ref{fig:pz_dependence}, the  $Z_{mp}$ and $Z_{\slashed{p}}-Z_{mp}$ at fixed $z\sim 0.5$ fm are plotted as a function of $p^R_z/\mu_R=1/\sqrt{r}$. From this figure, one can see   that the NPR factors, both real   and imaginary parts,  only show the dependence on  $r$ regardless of the values of $p_z^R$ or $\mu_R$, with  $p_z^R/\mu_R<0.4$.


\begin{figure}[tbp]
\includegraphics[width=.48\textwidth]{images/Z_pt_dep.pdf}
\caption{The  inverse of the minimum projection renormalization factor $1/Z_{mp}$ and the difference $1/Z_{\slashed{p}}-1/Z_{mp}$, as a function of $z$ with the same $p^R_z$ and $\mu_R$, but different $p^R_t$.  The crosses correspond to $p_t^R=0.9$ GeV and the open boxes (circles) are for $p_t^R=2.7$ GeV. Most results show mild dependence on $p^R_t$.} \label{fig:pt_dependence}
\end{figure}



In Fig.~\ref{fig:pt_dependence},  we show $1/Z_{mp}$ and $1/Z_{\slashed{p}}-1/Z_{mp}$ as functions of the Wilson link length $z$, with the same $\mu_R$=3.2 GeV and $p_z^R=1.4$ GeV and two different values of $p_t^R$=0.9 GeV and 2.7 GeV.  As shown in this figure, the  $1/Z_{mp}$ and $1/Z_{\slashed{p}}-1/Z_{mp}$ with the two different $p_t^R$'s are close  to each other for all $z$ (the curves with the same color).    This  is consistent with the 1-loop matching formula.   At $z<0.3$ fm, the real part of $1/Z_{\slashed{p}}-1/Z_{mp}$  with the two different $p_t^R$'s can be slightly nonzero, but it is  still  smaller than $1/Z_{\slashed{p}}$ by  two orders of magnitude.


In the following  we will take $p_z^R$ to be zero and estimate the systematic uncertainty from the $p_z^R$ dependence by varying the $p_z^R$.


In the calculation of nucleon matrix element, we use Gaussian momentum smearing~\cite{Bali:2016lva} for the quark field
\begin{align}\label{eq:moms}
\psi(x)\to &S_\text{mom}\psi(x) = \frac{1}{1+6\alpha}\nonumber\\
&\left[\psi(x) + \alpha \sum_j U^{APE}_j(x)e^{ik\hat{e}_j}\psi(x+\hat{e}_j)\right]\, ,
\end{align}
where $k$ is the desired momentum,  {$U^{APE}_j(x)$ are the APE smeared} gauge links in the $j$ direction, and $\alpha$ is a tunable parameter as in traditional Gaussian smearing.

Such a momentum source is designed to increase the overlap with nucleons of the desired boost momentum and we are able to reach
higher-boosted momentum for the nucleon states than
in the  previous work~\cite{Chen:2017mzz}. Although in the  exploratory
study, we varied the  Gaussian smearing radius to better
overlap with the largest momentum used in the calculation, the field smearing  is still centered around
zero momentum in momentum space. When we switch
to the momentum smearing, the  smearing center will be shifted to momentum $O(k)$, which will immediately allow us to reach higher boost momenta with better signal-to-noise ratios in the matrix elements. 
In this work, we use  two  values of nucleon boost momenta, $P_z=n \frac{2\pi}{L}$, with $n \in \{4,5\}$, which corresponds to 1.8 and 2.3~GeV.



\begin{figure*}[htbp]
\includegraphics[width=.49\textwidth,height= 0.35\textwidth]{images/quasipdf-R15-p4.pdf}
\includegraphics[width=.49\textwidth,height= 0.35\textwidth]{images/quasipdf-I15-p4.pdf}\\
\includegraphics[width=.49\textwidth,height= 0.35\textwidth]{images/quasipdf-R15-p5.pdf}
\includegraphics[width=.49\textwidth,height= 0.35\textwidth]{images/quasipdf-I15-p5.pdf}
\caption{The real (left) and imaginary (right) parts of the isovector nucleon matrix elements for unpolarized PDFs as functions of $z$ at different momenta, with $P_z=\frac{8\pi}{L}$=1.8 GeV (top) and $\frac{10\pi}{L}=$2.3~GeV (bottom) respectively. The RI/MOM renormalization factors with $\{\mu, p_z^R\}=\{3.2, 0\}$ GeV and   the normalization at $z=0$ are applied on the bare matrix elements to improve the visibility at large $z$.  At a given positive $z$ value, the data is slightly offset to show the ground-state matrix element from the fits using different ranges; from left to right they are: $t_\text{seq}\in$ [7,9], [7,10], [7,11], [8,11], and [9,11]. Different analyses are consistent within statistical errors while the fits with separation 7 and 8 have smaller  uncertainties compared  to   other cases.
}
\label{fig:bareME-tsep}
\end{figure*}




On the Lattice, we   calculate the time-independent and nonlocal  in $z$ direction  correlators of a nucleon with finite-$P_z$ boost
\begin{align}
\label{eq:qlat}
\widetilde{h}_\text{lat}(z,P_z,{\Gamma};a^{-1}) =
  \left\langle 0;\vec{P} \right|O_\Gamma(z)
  \left| 0;\vec{P} \right\rangle.
\end{align}
 {Here the state $|0; \vec{P}\rangle$ represents the ground (nucleon) state with momentum $\vec{P}=\{0,0,P_z\}$}.  $\Gamma=\gamma^t$ is used for the unpolarized parton distribution.



\begin{figure*}[htbp]
\includegraphics[width=.49\textwidth,height= 0.35\textwidth]{images/3ptV2ptZ12tFit7R-jack-comp-p4.pdf}
\includegraphics[width=.49\textwidth,height= 0.35\textwidth]{images/3ptV2ptZ12tFit7I-jack-comp-p4.pdf}\\
\includegraphics[width=.49\textwidth,height= 0.35\textwidth]{images/3ptV2ptZ10tFit7R-jack-comp-p5.pdf}
\includegraphics[width=.49\textwidth,height= 0.35\textwidth]{images/3ptV2ptZ10tFit7I-jack-comp-p5.pdf}
\caption{The real (left) and imaginary (right) parts of the renormalized isovector nucleon matrix elements for unpolarized PDFs with ${P_z, z}={8\pi/L, 12}$ (top) and ${10\pi/L, 10}$ (bottom) which correspond to $zP_z\sim9.5$. The data points and the band predicted by the fit using $t_\text{seq}\in$ [7, 11] agree with each other well.}
\label{fig:bareME-tsep2}
\end{figure*}


\begin{figure}[tbp]
\includegraphics[width=0.48\textwidth,height= 0.3\textwidth]{images/comp-p45-16-Re.pdf}
\includegraphics[width=0.48\textwidth,height= 0.3\textwidth]{images/comp-p45-16-Im.pdf}
\caption{The renormalized quasi-PDF matrix elements with $P_z=$1.8 and 2.3 GeV, using the minimal projection with $p_z^R=0$ and $\mu_R=3.2$ GeV, as   function of $zP_z$.} \label{fig:renormalized_h-tilde}
\end{figure}

As  the nucleon boost momentum increases, one anticipates that excited-state contributions  are more severe; therefore, a careful study of the excited-state contamination is necessary.
To do so, we calculate the nucleon matrix element $\widetilde{h}_{\rm lat}$ at five source-sink separations $t_\text{seq}\in\{7,8,9,10,11\}\times 0.086$~fm, with $\{4,4,8,8,16\}$~measurements on each of 2005 gauge configurations respectively in the $P_z=1.8$ GeV case, and  of 2,000 configurations in the $P_z=2.3$ GeV case. We use a multigrid algorithm~\cite{Babich:2010qb,Osborn:2010mb} with the Chroma software package~\cite{Edwards:2004sx} to speed up the inversion of the quark propagator.
Following  Ref.~\cite{Bhattacharya:2013ehc}, each three-point (3pt) correlator $C_\Gamma^{(3\text{pt})} (t, t_\text{seq})$ can be decomposed as (assuming the source is at $t= 0$)
\begin{align} \label{eq:fit}
C^\text{3pt}(t,t_\text{seq}; P_z, {\Gamma}) &=
   {\cal A}_0^2 \langle 0 | O_\Gamma | 0 \rangle  e^{-E_0t_\text{seq}} \\
   &+{\cal A}_1^2 \langle 1 | O_\Gamma | 1 \rangle  e^{-E_1t_\text{seq}} \nonumber\\
   &+{\cal A}_1{\cal A}_0 \langle 1 |O_\Gamma | 0 \rangle  e^{-E_1 (t_\text{seq}-t)} e^{-E_0 t} \nonumber\\
   &+{\cal A}_0{\cal A}_1 \langle 0 | O_\Gamma | 1 \rangle  e^{-E_0 (t_\text{seq}-t)} e^{-E_1 t} + \ldots \,,\nonumber
\end{align}
where  $|n\rangle$  with $n
> 0$ represents the excited states. The operator is inserted at time
$t$, and the nucleon state is annihilated at the sink time
$t_\text{seq}$( which is also the source-sink separation).
The spectrum weights ${\cal A}_{0,1}$ and energies $E_{0,1}$ in Eq.~(\ref{eq:fit}) can be obtained from the two-point (2pt) correlator: 
\begin{align} \label{eq:fit2}
C^\text{2pt}(t_\text{seq}; P_z) &=
   {\cal A}_0^2 e^{-E_0t_\text{seq}} +{\cal A}_1^2 e^{-E_1t_\text{seq}} + \ldots\ .
\end{align}





Eventually we apply  the joint fit with the 3pt functions at several $t_\text{seq}$ and   2pt function using the following form~\cite{Bhattacharya:2013ehc}:
\begin{align}
&\frac{C^\text{3pt}(t,t_\text{seq})}{C^\text{2pt}(t_\text{seq})}= \nn\\
&=\frac{\tilde{h}_{\text{lat}}+C_2(e^{-\Delta Et}+e^{-\Delta E(t_\text{seq}-t)})+C_3e^{-\Delta E t_\text{seq}}}{1+C_1 e^{-\Delta Et_\text{seq}}},\nonumber\\
&C^\text{2pt}(t)=C_0e^{-E_0t}(1+C_1 e^{-\Delta Et}),
\end{align}
with $\Delta E= E_1-E_0$.  $C_{0,1,2,3}$ and $E_{0,1}$ are free parameters. We limit the range of $t$ as  $t\in[1,t_\text{seq}-1]$ for 3pt/2pt ratio and $t\in[7, 11]$ for the 2pt to make the $\chi^2/d.o.f.$ of the fit to be ${\cal O} (1)$.  Using the ratio of 3pt/2pt instead of the 3pt function itself can improve the stability of the fit, especially when $C^\text{2pt}(t)$ with $t<7$ is included in the fit.




In Fig.~\ref{fig:bareME-tsep}, we show the ground-state nucleon matrix elements $\tilde{h}_{lat}(z,P_z,\gamma_t)$ obtained from five fits: using the separations $t_\text{seq}\in$ [7, 9], [7, 10], [7,11], [8,11], and [9,11] (The data points correspond to the same $z$ but are shifted horizontally to enhance the visibility). The  data are further normalized by multiplying  the renormalization factor with $\{\mu_R, p_z^R\}=\{3.2, 0\}$ GeV and the real part normalized to 1 at $z=0$. From this figure, one can see that  there is no clear signal for excited-state contributions in any of these analyses.    If the data with    smallest two separations are dropped, uncertainties are getting much larger.     In the fit,  we keep the $C_3$ term to make a moderate estimate of the uncertainty even when this term is not statistically significant.


For  a comparison between data and the fit, we show our results at large $z$ like  $({P_z, z})=({8\pi/L, 12a})$ and $({10\pi/L, 10a})$ with  $t_\text{seq}\in$ [7,11] in Fig.~\ref{fig:bareME-tsep2}. In these  spatial separations,  the real part of  matrix element seems to be negative.   The ground-state contribution  obtained from the fit is shown as the black band. As one can see, most data can be well described in the fit and thereby we use the two-state fits and the interval  $t_\text{seq}\in$ [7, 11] to obtain the results   in the rest part of this paper.



The renormalized quasi-PDF matrix elements with two values of $P_z$ are plotted in Fig.~\ref{fig:renormalized_h-tilde}, as  function of $zP_z$ for  $p^R_z=0$ and $\mu_R=3.2$ GeV. The results with different $P_z$ are consistent with each other within statistical uncertainties.  This  indicates that power corrections due to   higher-twist effects might not be sizable.


\subsection{Systematic  {Uncertainties}}


\begin{figure}
\includegraphics[width=0.48\textwidth]{images/error_pzR=0.pdf}
\includegraphics[width=0.48\textwidth]{images/error_pzR=3.pdf}
\caption{\label{fig:error} Different contributions to the systematic errors. See the text for detailed information.  }
\end{figure}


In this subsection, we will consider four systematic uncertainties from: Fourier transformation (FT),  unphysical scales $p_z^R$ and $\mu_R$, projection used in the RI/MOM scheme, and inversion of the matching coefficient. 



\begin{figure}[tbp]
\includegraphics[width=.48\textwidth]{images/matching_cutoff.pdf}
\includegraphics[width=.48\textwidth]{images/matching_derivative.pdf}
\caption{\label{fig:matchingPDF} The quasi-PDF (dashed-red) with nucleon boost momentum 2.3 GeV and matched PDF (solid-black) at $\mu=2$~GeV using minimal projection, with the RI/MOM parameters $p_z^R=0$ and $\mu_R=3.2$~GeV. The upper and lower figures are obtained  using the  derivative and cutoff methods to perform  Fourier transformation. The matching strategy has an important impact on the final results for PDFs.}
\end{figure}


In the following, we explain the details to include these systematic uncertainties.

 {1) Fourier transformation.  As shown in Fig.~\ref{fig:renormalized_h-tilde},  the  $\widetilde{h}_R(z)$ with $P_z$=2.3 GeV is consistent with zero when $z>12a$. Thus in the standard matching from quasi-PDF to PDF, it is reasonable to truncate the results at $z=12a$. With this spirit,  the quasi-PDF and matched PDF using the standard FT are shown in the upper panel of Fig.~\ref{fig:matchingPDF}, from which one can see  the matched PDF shows an oscillatory behavior. A ``derivative'' method was proposed in Ref. ~\cite{Lin:2017ani}  to cure  this  oscillatory behavior. To be concrete, one takes the derivative of the renormalized nucleon matrix elements $\partial_z\widetilde{h}_R(z)$, whose Fourier transform differs from the original matrix element in a known way:
%
\begin{equation}
\label{eq:derivative}
\widetilde{q}_R(x) = \int^\infty_{-\infty} \frac{dz}{2\pi} \frac{ie^{i x P_z z}}{x} \partial_z\widetilde{h}_R(z),
\end{equation}
%
provided that $\widetilde{h}_R(z)$ goes to zero as $|z|\to\infty$. With the same truncation, the result is shown in the lower panel of Fig.~\ref{fig:matchingPDF} and apparently  the oscillatory behavior is less severe. Besides, results obtained using the derivative method is consistent with the standard FT method in most kinematics region except   at  small $x$. This is anticipated  as two methods only differ at the large $z$ region where we have made the truncation.  We show  the  difference as the dot-dashed-blue line in Fig.~\ref{fig:error}, together with the error from varying the truncation from $z=10a$ to $14a$ (dotted-green line). }

\begin{figure}[tbp]
\includegraphics[width=0.45\textwidth]{images/matching_error_pzR=0.pdf}
\includegraphics[width=0.45\textwidth]{images/matching_error_pzR=3.pdf}
\caption{\label{fig:mathcing_error_test} Effects of inversion matching formula using minimal projection: The solid-black, dotted-red, dotted-blue, and dot-dashed-green lines represent CT14nnlo PDF, applying inverse matching from CT14nnlo PDF~\cite{Dulat:2015mca} to quasi-PDF, applying matching again to get back to the PDF, the difference between PDF with iterative matching and the original CT14nnlo PDF. The upper (lower) figure corresponds to $p_z^R=0$ (1.4 GeV). These plots show that the method we used to invert the matching formula is less reliable for small $|x|$. The difference shown by the dot-dashed-green curve is taken into account into our systematic error.}
\end{figure}

 {2) Unphysical scales $p_z^R$ and $\mu_R$. There are two unphysical scales $p_z^R$ and $\mu_R$ introduced in RI/MOM.  In principle, when matching the quasi-PDF matrix element onto lightcone PDF, the dependence on these two scales  in the matrix element   should exactly  cancel with that in the matching kernel.  However, since the quasi-PDF matrix element is non-perturbatively renormalized  on the Lattice, while the matching coefficient is  calculated at one-loop order in perturbation theory, there will be residual dependence on these two scales after the perturbative matching. To estimate the residual $p_z^R$ and $\mu_R$ dependence, we choose   {$p_z^R=0$ GeV and $\mu_R=3.2$ GeV} as the central value,  and vary $p_z^R$ from -1.4 to 1.4 GeV (dashed-red line) and $\mu_R$ from 2.4 to 3.9 GeV (thick-solid-orange line). The   difference between  these matched PDFs is treated  as the systematics of the residual dependence on unphysical scales, in Fig.~\ref{fig:error}. As shown in the figure, the systematic uncertainty due to the $\mu_R$ dependence is small compared  to the other sources, but the residual $p_z^R$ dependence could be sizable.}

 {3) Dependencies on the projection. There are two projections discussed in this work:  the minimum and ``$\slashed p$" projections.  With $p_z^R=0$, the projection dependence in both the NPR factor are less than 5\%  for all the $z$, and vanishes in the 1-loop perturbative matching.  Thus one can expect that the difference due to the projections is also small, as depicted  in the long-dashed-magenta lines in Fig.~\ref{fig:error}.}




4) Inversion of matching. To extract the PDF from the quasi-PDF, one needs to invert the factorization formula Eq. (\ref{eq:fact}). This could be done by   changing the sign of $\alpha_s$ in $C$,  and convoluting the new matching coefficient with $\widetilde{q}$. More explicitly, we  have
\begin{align}\label{eq:inverse_fact}
q(x,\mu)=&\int_{-\infty}^\infty {dy\over |y|}\: C'\left({x\over y},r,\frac{yP_z}{\mu},\frac{yP_z}{p_z^R}\right)\widetilde{q}(y,P_z, p^R_z,\mu_R)\nonumber\\
&+\mathcal{O}\left({M^2\over P_z^2},{\Lambda_{\text{QCD}}^2\over x^2 P_z^2},\alpha_s^2\right)\,,
\end{align}
where $C'=C(\alpha_s\to-\alpha_s)$. We estimate the error due to inverting the factorization formula  by  {starting from the PDF from a global analysis~\cite{Dulat:2015mca}, applying Eq.(\ref{eq:fact}) and then Eq. (\ref{eq:inverse_fact}) to return to the PDF. This manipulation should give the same lightcone PDF}. However, since the matching  is   only accurate up to $O(\alpha_s)$, the two results would  differ, and the difference gives a good estimate of the systematic error coming from the inversion and higher order corrections. This is shown in Fig.~\ref{fig:mathcing_error_test}, from which one can find   that the error only  becomes sizable  when $|x|$ is small.  {This is expected because the relevant momentum scale is $xP^z$ such that higher order corrections become large at small $x$}.

There are more sophisticated methods to invert the factorization, such as using a recursion procedure. However, as we can see in Fig.~\ref{fig:error}, the systematic error caused by the matching procedure (thick-dashed-cyan line) is also smaller than those from the first two sources in most regions.


 As shown in Fig.~\ref{fig:error}, one can find  that the dominant   uncertainties arise  from  the $p_z^R$ dependence and  from FT. With $p_z^R=0$, the uncertainties  from different  projections, and  inversion and the matching are typically  less than $10\%$ except in the region with very small $|x|$. The uncertainty  from the $\mu_R$ dependence is even smaller. With  $p_z^R=1.4$ GeV, uncertainties from these  three sources are getting  larger in magnitude, but still smaller than those from the two major sources. 

